\chapter{Ray Tracing}

O \textit{Ray Tracing} (Traçado de Raios) representa uma mudança fundamental de paradigma na síntese de imagens digitais. Enquanto a rasterização projeta geometria 3D em um plano 2D de forma linear e local, o Ray Tracing simula o comportamento físico da luz ao seguir trajetórias de raios através da cena. Como observa \textcite{shirley2021fundamentals}, esta técnica é inerentemente recursiva e global, permitindo a simulação natural de fenômenos ópticos complexos como sombras suaves, reflexões especulares, refrações e iluminação indireta. Com o advento de hardware especializado em GPUs modernas, o Ray Tracing deixou de ser restrito ao render offline de alta qualidade e passou a integrar pipelines de tempo real em jogos e simulações profissionais de alta fidelidade.

\section{Evolução Histórica e Modelos de Renderização}

A trajetória das técnicas de traçado de raios reflete o aumento exponencial do poder computacional e a sofisticação dos modelos físicos de transporte de luz. Historicamente, o desenvolvimento seguiu os seguintes marcos:

\begin{itemize}
    \item \textbf{Ray Casting (Appel, 1968):} O ponto de partida da visibilidade 3D. O algoritmo lança um único raio por pixel para determinar qual objeto é visível para o observador. Embora resolva o problema de superfícies ocultas, ele não lida com reflexões, refrações ou sombras de forma recursiva, limitando-se a modelos de iluminação local.
    \item \textbf{Recursive Ray Tracing (Whitted, 1980):} Turner Whitted revolucionou a área ao introduzir o transporte recursivo de luz. Ao atingir uma superfície, o raio primário pode gerar novos raios de reflexão e refração ideais (espelhadas), além de raios de sombra direcionados às fontes de luz pontuais. Este modelo permitiu, pela primeira vez, a renderização de objetos transparentes e reflexivos de forma convincente.
    \item \textbf{Distributed Ray Tracing (Cook et al., 1984):} Introduziu a amostragem estocástica (Monte Carlo) para simular efeitos de "foco suave". Ao lançar múltiplos raios levemente perturbados por pixel, o algoritmo permitiu renderizar reflexões borradas (\textit{glossy}), sombras suaves de luzes de área, \textit{motion blur} e profundidade de campo (\textit{depth of field}), aproximando-se da fotografia real.
    \item \textbf{Path Tracing (Kajiya, 1986):} A técnica definitiva que utiliza integração de Monte Carlo para resolver a Equação de Renderização completa. O Path Tracing permite a simulação de luz indireta difusa (\textit{color bleeding}), cáusticas e outros fenômenos complexos de iluminação global. Atualmente, é o padrão ouro na indústria cinematográfica e de efeitos visuais.
\end{itemize}

\section{Fundamentos de Radiometria}

Para construir um renderizador fisicamente correto (\textit{Physically Based Rendering} - PBR), é essencial compreender as grandezas físicas que regem o transporte de energia luminosa:

\begin{itemize}
    \item \textbf{Fluxo Radiante ($\Phi$):} A potência total emitida, refletida ou recebida por uma superfície, medida em Watts ($W$).
    \item \textbf{Irradiância ($E$):} O fluxo recebido por unidade de área de uma superfície ($W/m^2$). É fundamental para o cálculo de iluminação incidente.
    \item \textbf{Radiância ($L$):} A grandeza mais importante em Ray Tracing. Representa o fluxo por unidade de área por unidade de ângulo sólido ($W/(m^2 \cdot sr)$). A radiância é invariante ao longo de um raio no vácuo, o que justifica o lançamento de raios a partir da câmera para "buscar" a cor na cena.
    \item \textbf{Ângulo Sólido ($\omega$):} A extensão tridimensional de um ângulo plano, medida em esferorradianos (sr). Uma esfera completa engloba $4\pi$ sr.
\end{itemize}

\section{A Geometria Analítica do Raio}

Um raio é uma estrutura geométrica que representa a trajetória retilínea da luz. Formalmente, é uma semirreta definida por um ponto de origem $\mathbf{O}$ e um vetor de direção unitário $\mathbf{D}$.

\begin{formula}{Equação Paramétrica do Raio}
    R(t) = \mathbf{O} + t\mathbf{D}, \quad t \in [t_{min}, t_{max}]
\end{formula}

O parâmetro escalar $t$ representa a distância percorrida a partir da origem. Para que o raio avance na direção de interesse, devemos garantir que $t > 0$. Em implementações práticas, define-se um limite inferior $t_{min} \approx 0.0001$ para mitigar artefatos de precisão numérica (\textit{shadow acne}) e um limite superior $t_{max}$ para limitar o alcance da busca espacial.

\section{Sistemas de Coordenadas e Câmera}

A geração de raios primários exige a transformação de coordenadas do espaço da imagem (2D) para o espaço do mundo (3D). Este processo simula o funcionamento de uma câmera \textit{pinhole}.

\subsection{Câmera Pinhole e Viewport}

Considere uma câmera na origem $\mathbf{E}$ (olho), olhando para a direção $-\mathbf{W}$ (vetor unitário), com vetor "up" $\mathbf{V}$ e vetor lateral $\mathbf{U}$. Definimos um plano de imagem a uma distância $d$ do olho. Para um pixel $(i, j)$ em uma imagem de resolução $N_x \times N_y$:
\begin{align}
    u = l + (r-l)\frac{i+0.5}{N_x} \\
    v = b + (t-b)\frac{j+0.5}{N_y}
\end{align}
Onde $l, r, b, t$ são as coordenadas das bordas esquerda, direita, baixo e topo do viewport. A direção do raio primário $\mathbf{D}$ é dada por:
\begin{equation}
    \mathbf{D} = -d\mathbf{W} + u\mathbf{U} + v\mathbf{V}
\end{equation}

\subsection{Lentes Reais e Profundidade de Campo}

Para simular o efeito de desfoque (\textit{Depth of Field}), substituímos o ponto matemático da câmera por um disco de abertura. Em vez de todos os raios partirem da mesma origem, amostramos um ponto aleatório na lente e calculamos a direção tal que o raio passe pelo ponto de foco no plano focal. Isso resulta em objetos fora do plano focal aparecendo borrados (\textit{Circle of Confusion}).

\section{Interseção Raio-Geometria}

\section{Sistemas de Coordenadas e Câmera}

A geração de raios primários exige a transformação de coordenadas do espaço da imagem (2D) para o espaço do mundo (3D). Este processo simula o funcionamento de uma câmera \textit{pinhole}.

\subsection{Câmera Pinhole e Viewport}

Considere uma câmera na origem $\mathbf{E}$ (olho), olhando para a direção $-\mathbf{W}$ (vetor unitário), com vetor "up" $\mathbf{V}$ e vetor lateral $\mathbf{U}$. Definimos um plano de imagem a uma distância $d$ do olho. Para um pixel $(i, j)$ em uma imagem de resolução $N_x \times N_y$:
\begin{align}
    u = l + (r-l)\frac{i+0.5}{N_x} \\
    v = b + (t-b)\frac{j+0.5}{N_y}
\end{align}
Onde $l, r, b, t$ são as coordenadas das bordas esquerda, direita, baixo e topo do viewport. A direção do raio primário $\mathbf{D}$ é dada por:
\begin{equation}
    \mathbf{D} = -d\mathbf{W} + u\mathbf{U} + v\mathbf{V}
\end{equation}

\subsection{Lentes Reais e Profundidade de Campo}

Para simular o efeito de desfoque (\textit{Depth of Field}), substituímos o ponto matemático da câmera por um disco de abertura. Em vez de todos os raios partirem da mesma origem, amostramos um ponto aleatório na lente e calculamos a direção tal que o raio passe pelo ponto de foco no plano focal. Isso resulta em objetos fora do plano focal aparecendo borrados (\textit{Circle of Confusion}).

\section{Interseção Raio-Geometria}

O cálculo de interseção é a operação mais custosa de um renderizador, ocupando frequentemente mais de 90\% do tempo de CPU/GPU.

\subsection{Interseção com Esferas: Derivação Algébrica Detalhada}

Uma esfera é definida pelo centro $\mathbf{C}$ e raio $r$. Qualquer ponto $\mathbf{P}$ na superfície satisfaz $\|\mathbf{P}-\mathbf{C}\|^2 = r^2$. Ao substituir a equação do raio $R(t) = \mathbf{O} + t\mathbf{D}$, obtemos:
\begin{equation}
    (\mathbf{O} + t\mathbf{D} - \mathbf{C}) \cdot (\mathbf{O} + t\mathbf{D} - \mathbf{C}) - r^2 = 0
\end{equation}
Expandindo os termos e organizando como uma equação quadrática na forma $at^2 + bt + c = 0$:
\begin{itemize}
    \item $a = \mathbf{D} \cdot \mathbf{D} = 1$ (assumindo direção normalizada)
    \item $b = 2\mathbf{D} \cdot (\mathbf{O} - \mathbf{C})$
    \item $c = (\mathbf{O} - \mathbf{C}) \cdot (\mathbf{O} - \mathbf{C}) - r^2$
\end{itemize}
O discriminante $\Delta = b^2 - 4ac$ determina o número de interseções. Se $\Delta < 0$, o raio erra a esfera. Se $\Delta = 0$, o raio é tangente. Se $\Delta > 0$, existem duas interseções; em Ray Tracing, estamos interessados no menor valor de $t$ que seja maior que $t_{min}$.

\subsection{Interseção com Triângulos: O Algoritmo de Möller-Trumbore}

O algoritmo de Möller-Trumbore é amplamente utilizado por ser rápido e não exigir o armazenamento das normais dos planos das faces. Ele utiliza coordenadas baricêntricas para verificar se o ponto de interseção com o plano do triângulo (vértices $A, B, C$) está dentro dos limites da geometria.

Definimos as arestas $\mathbf{E}_1 = B - A$ e $\mathbf{E}_2 = C - A$. Um ponto no triângulo é expresso como $P(u, v) = A + u\mathbf{E}_1 + v\mathbf{E}_2$. Igualando ao raio:
\begin{equation}
    \mathbf{O} + t\mathbf{D} = A + u\mathbf{E}_1 + v\mathbf{E}_2 \Rightarrow [-D, \mathbf{E}_1, \mathbf{E}_2] \begin{bmatrix} t \\ u \\ v \end{bmatrix} = \mathbf{O} - A
\end{equation}
Aplicando a Regra de Cramer, obtemos $t, u, v$. A interseção é válida se e somente se:
\begin{formula}{Condições de Interseção em Triângulos}
    t > t_{min}, \quad u \ge 0, \quad v \ge 0, \quad (u + v) \le 1
\end{formula}

\subsection{Interseção com AABB: O Método das Fatias (Slabs)}

Volumes envolventes alinhados aos eixos (AABBs) são definidos por dois pontos: $\mathbf{P}_{min}$ e $\mathbf{P}_{max}$. O método de Kay-Kajiya trata o AABB como a interseção de três "fatias" infinitas (uma para cada eixo).

Para cada eixo $i \in \{x, y, z\}$:
\begin{align}
    t_{1} = (P_{min, i} - O_i) / D_i \\
    t_{2} = (P_{max, i} - O_i) / D_i \\
    t_{near, i} = \min(t_1, t_2), \quad t_{far, i} = \max(t_1, t_2)
\end{align}
O intervalo de interseção global é dado por:
\begin{equation}
    t_{enter} = \max(t_{near, x}, t_{near, y}, t_{near, z}), \quad t_{exit} = \min(t_{far, x}, t_{far, y}, t_{far, z})
\end{equation}
O raio atinge o volume se $t_{enter} \le t_{exit}$ e $t_{exit} > 0$. Esta operação é extremamente eficiente pois utiliza apenas subtrações e multiplicações (o inverso da direção $1/D_i$ pode ser pré-calculado).

\section{Estruturas de Aceleração Espacial: BVH e KD-Trees}

Em cenas de produção com milhões de triângulos, realizar testes lineares contra toda a geometria é computacionalmente proibitivo ($O(N)$). Estruturas de dados espaciais reduzem essa complexidade para $O(\log N)$.

\subsection{Bounding Volume Hierarchy (BVH)}

A BVH é a estrutura mais comum em Ray Tracing moderno (incluindo hardware de RT). Ela organiza os objetos em uma árvore onde cada nó armazena um AABB que engloba todos os seus descendentes.
\begin{itemize}
    \item \textbf{Construção:} O conjunto de primitivas é dividido recursivamente. A escolha do plano de corte é crítica para a performance.
    \item \textbf{Surface Area Heuristic (SAH):} Em vez de dividir o espaço ao meio, a SAH estima o custo de interseção baseado na probabilidade de um raio atingir um nó, que é proporcional à sua área de superfície. O custo de um corte é $C = C_{trav} + P(A) \cdot N_A \cdot C_{int} + P(B) \cdot N_B \cdot C_{int}$, onde $P(X) = SurfaceArea(X) / SurfaceArea(Parent)$.
    \item \textbf{Travessia:} Um raio é testado contra o AABB da raiz. Se atingir, o processo se repete para os filhos. Se um nó folha é alcançado, o raio é testado contra as primitivas reais.
\end{itemize}

\subsection{Grades Espaciais e KD-Trees}

Enquanto a BVH subdivide os objetos, as KD-Trees subdividem o espaço. Grades uniformes são excelentes para cenas com distribuição homogênea de objetos, mas sofrem com o problema do "estádio no deserto" (grandes áreas vazias consumindo memória). KD-Trees resolvem isso com subdivisões adaptativas, sendo historicamente populares para busca de fótons em \textit{Photon Mapping}.


\section{A Equação de Renderização: O Modelo Unificado}

Formulada independentemente por \textcite{kajiya1986rendering} e Immel et al. em 1986, a Equação de Renderização é a base teórica de toda a síntese de imagem baseada em física. Ela descreve o balanço de energia luminosa (radiância) em um sistema fechado, assumindo que a luz viaja instantaneamente e que o meio não interfere na propagação (vácuo).

A forma integral da equação para a radiância de saída $L_o$ em um ponto $p$ na direção $\omega_o$ é:
\begin{formula}{A Equação de Renderização}
    L_o(p, \omega_o) = L_e(p, \omega_o) + \int_{\Omega} f_r(p, \omega_i, \omega_o) L_i(p, \omega_i) (\omega_i \cdot \mathbf{n}) d\omega_i
\end{formula}

Componentes da equação:
\begin{itemize}
    \item $L_o(p, \omega_o)$: A radiância total saindo de $p$ na direção $\omega_o$. Este é o valor que a câmera eventualmente captura.
    \item $L_e(p, \omega_o)$: Radiância emitida diretamente pelo ponto (se ele for parte de uma fonte de luz).
    \item $\int_{\Omega} \dots d\omega_i$: Integral sobre o hemisfério $\Omega$ centrado na normal $\mathbf{n}$. Ela soma a luz vinda de todas as direções incidentes $\omega_i$.
    \item $f_r(p, \omega_i, \omega_o)$: A BRDF (\textit{Bidirectional Reflectance Distribution Function}). Define a proporção de luz vinda de $\omega_i$ que é refletida para $\omega_o$. Deve respeitar a conservação de energia e a reciprocidade de Helmholtz.
    \item $L_i(p, \omega_i)$: Radiância incidente vinda de outras partes da cena. Em um sistema de Ray Tracing, este termo é resolvido recursivamente: $L_i(p, \omega_i) = L_o(p', -\omega_i)$, onde $p'$ é o primeiro ponto atingido pelo raio disparado de $p$ na direção $\omega_i$.
    \item $(\omega_i \cdot \mathbf{n})$ ou $\cos \theta_i$: O termo de atenuação de Lambert. Representa a redução da densidade de fluxo quando a luz atinge a superfície em um ângulo oblíquo.
\end{itemize}

Resolver esta integral é o desafio central do Ray Tracing. Métodos como o \textit{Path Tracing} utilizam integração de Monte Carlo para estimar o valor da integral amostrando o hemisfério de forma estocástica.

\section{Raios Secundários e Fenômenos Ópticos}

A versatilidade do Ray Tracing reside na sua capacidade de simular efeitos globais simplesmente disparando raios adicionais a partir do ponto de impacto primário.

\subsection{Sombra e Oclusão}

Diferente da rasterização (que usa \textit{shadow maps} ou volumes), o Ray Tracing resolve sombras disparando raios de oclusão (\textit{shadow rays}) em direção às luzes. Se um raio disparado de $p$ em direção à luz $L$ encontrar qualquer objeto com $t < \|L-p\|$, o ponto está na sombra. Luzes de área (não pontuais) produzem sombras suaves naturalmente se dispararmos múltiplos raios para pontos aleatórios na superfície da lâmpada, simulando a \textbf{penumbra}.

\subsection{Reflexão e Refração Especular}

Materiais perfeitamente polidos (espelhos) ou dielétricos (vidro) geram raios secundários em direções determinísticas:
\begin{itemize}
    \item \textbf{Reflexão:} A direção refletida $\mathbf{r}$ para um raio incidente com direção $\mathbf{d}$ e normal $\mathbf{n}$ é $\mathbf{r} = \mathbf{d} - 2(\mathbf{d}\cdot\mathbf{n})\mathbf{n}$.
    \item \textbf{Refração:} Regida pela Lei de Snell ($n_1 \sin \theta_1 = n_2 \sin \theta_2$). O raio muda de direção ao passar de um meio para outro com índices de refração (\textit{IOR}) diferentes.
    \item \textbf{Efeito Fresnel:} Em materiais reais, a razão entre luz refletida e refratada depende do ângulo de incidência. Observadores próximos à perpendicular de uma superfície de vidro veem mais refração; em ângulos rasos (\textit{grazing angles}), a reflexão domina.
\end{itemize}

\section{Path Tracing e Monte Carlo}

O Path Tracing resolve a integral da Equação de Renderização através da amostragem estocástica de caminhos de luz. Em vez de calcular a integral analiticamente (o que é impossível para a maioria das cenas), estimamos o valor médio da radiância através de $N$ amostras aleatórias no domínio do hemisfério $\Omega$.

\begin{formula}{Estimador de Monte Carlo para Radiância}
    L_o(p, \omega_o) \approx \frac{1}{N} \sum_{k=1}^N \frac{f_r(p, \omega_{i,k}, \omega_o) L_i(p, \omega_{i,k}) \cos \theta_{i,k}}{p(\omega_{i,k})}
\end{formula}

Onde $p(\omega_{i,k})$ é a função de densidade de probabilidade (PDF) da amostra $k$. A precisão do resultado aumenta conforme $1/\sqrt{N}$, o que significa que para reduzir o ruído visível pela metade, precisamos quadruplicar o número de raios por pixel.

\subsection{Técnicas de Redução de Variância}

A variância no Path Tracing manifesta-se visualmente como ruído (\textit{fireflies} ou granulação). Para acelerar a convergência sem viés (\textit{bias}), utilizamos:

\begin{itemize}
    \item \textbf{Importance Sampling:} Em vez de escolher direções $\omega_i$ uniformemente, amostramos direções que provavelmente trarão mais energia (conforme a BRDF ou apontando para luzes). Se a PDF for proporcional ao integrando, a variância é minimizada.
    \item \textbf{Russian Roulette:} Caminhos de luz podem, teoricamente, ter comprimento infinito. Para evitar recursão infinita e desperdício em caminhos de baixa contribuição, usamos a Roleta Russa. Dado uma probabilidade de sobrevivência $q$:
    \begin{equation}
        L_{new} = \begin{cases} L/q & \text{se } \xi < q \\ 0 & \text{se } \xi \ge q \end{cases}
    \end{equation}
    Onde $\xi$ é um número aleatório uniforme em $[0, 1)$. O termo $1/q$ garante que o estimador permaneça não-viesado.
    \item \textbf{Next Event Estimation (NEE):} Separa a iluminação direta da indireta. Em cada interseção, lança-se explicitamente um raio para um ponto aleatório em cada fonte de luz (\textit{shadow ray}). Isso resolve o problema de luzes pequenas que seriam raramente atingidas por amostragem aleatória do hemisfério.
    \item \textbf{Multiple Importance Sampling (MIS):} Introduzido por Veach, o MIS combina múltiplas estratégias de amostragem (ex: amostrar a BRDF e amostrar as luzes) ponderando-as de modo a minimizar a variância total, mesmo quando uma das estratégias falha (como uma luz muito grande ou um material muito especular).
\end{itemize}

\section{Modelos de Microfacetas e PBR}

A renderização baseada em física (\textit{Physically Based Rendering} - PBR) revolucionou a indústria ao fornecer um modelo matemático unificado que garante a consistência da iluminação sob diferentes condições.

\subsection{A BRDF de Cook-Torrance}

O modelo assume que superfícies macroscópicas são formadas por microfacetas microscópicas orientadas de forma estatística. A BRDF especular é composta por três termos principais:
\begin{equation}
    f_{spec} = \frac{D \cdot G \cdot F}{4(\omega_i \cdot \mathbf{n})(\omega_o \cdot \mathbf{n})}
\end{equation}

\begin{itemize}
    \item \textbf{D (Distribuição):} A função NDF (\textit{Normal Distribution Function}) define a fração de microfacetas orientadas na direção do \textit{half-vector} $\mathbf{h} = \text{normalize}(\omega_i + \omega_o)$. O modelo \textbf{GGX} é o padrão moderno devido ao seu "glint" característico e transição suave entre o pico especular e a área difusa.
    \item \textbf{G (Geometria):} Modela a auto-oclusão. Em ângulos rasos, uma microfaceta pode bloquear a luz incidente ou refletida de outra microfaceta. Isso é essencial para evitar que materiais pareçam excessivamente brilhantes nas bordas.
    \item \textbf{F (Fresnel):} Descreve a refletância em função do ângulo. A aproximação de \textbf{Schlick} simplifica o cálculo: $R(\theta) = R_0 + (1 - R_0)(1 - \cos \theta)^5$. Onde $R_0$ é a refletância em incidência normal, que define se o material é metal ou dielétrico.
\end{itemize}

\subsection{Conservação de Energia}

Um princípio fundamental do PBR é que a energia refletida nunca pode exceder a energia incidente. No código, isso é implementado garantindo que a soma das componentes difusa ($k_d$) e especular ($k_s$) seja menor ou igual a 1, e que o termo de Fresnel seja usado para ponderar essas componentes dinamicamente.

\section{Pseudocódigo: Kernel de Path Tracing}

Abaixo, apresentamos a estrutura de um integrador de caminhos (\textit{path integrator}) robusto com NEE (\textit{Next Event Estimation}) e Roleta Russa:

\begin{lstlisting}[language=C++, caption=Kernel de Path Tracing Moderno]
Color Radiance(Ray ray, Scene scene, int depth) {
    if (depth > max_depth) return Color(0);

    Hit hit = scene.intersect(ray);
    if (!hit.occurred()) return scene.backgroundColor(ray.dir);

    Color L = 0;
    Material mat = hit.material;

    // Emissao direta (atingiu uma luz)
    L += mat.emission;

    // Next Event Estimation (Shadow Rays)
    for (Light light : scene.lights) {
        ShadowRay sRay = sampleLight(hit.p, light);
        if (!scene.occluded(sRay)) {
            L += mat.eval(hit.wo, sRay.wi) * light.intensity * 
                 dot(hit.n, sRay.wi) / sRay.pdf;
        }
    }

    // Amostragem da BRDF para o proximo rebote
    Sample bSample = mat.sample(hit.wo, hit.n);
    Ray nextRay(hit.p, bSample.wi);
    
    // Roleta Russa
    float q = max(0.05f, 1.0f - bSample.albedo);
    if (random() > q) {
        L += bSample.weight * Radiance(nextRay, scene, depth + 1) / (1-q);
    }

    return L;
}
\end{lstlisting}

\section{Materiais Participantes e Espalhamento}

Efeitos como fumaça, névoa e a aparência da pele humana exigem a simulação de luz viajando através de volumes.

\begin{itemize}
    \item \textbf{Absorption and Out-scattering:} A radiância diminui à medida que o raio viaja pelo volume, sendo absorvida ou desviada para outras direções.
    \item \textbf{Phase Functions:} Análogo à BRDF para superfícies, a função de fase (ex: Henyey-Greenstein) define a distribuição angular do espalhamento dentro do volume.
    \item \textbf{Subsurface Scattering (SSS):} Ocorre quando a luz entra em um objeto translúcido, sofre múltiplos espalhamentos internos e sai em um ponto diferente. Essencial para renderizar materiais orgânicos com realismo fotográfico.
\end{itemize}

\section{Técnicas de Monte Carlo e Redução de Variância}

A eficiência de um renderizador de caminhos (\textit{Path Tracer}) depende de quão rápido ele consegue reduzir a variância (ruído) da imagem.

\subsection{Importância da Amostragem (Importance Sampling)}

Em vez de escolher direções $\omega_i$ uniformemente no hemisfério, escolhemos direções que possuem maior probabilidade de contribuir para a integral da Equação de Renderização. Se conhecemos a forma da BRDF ou a posição das luzes, podemos concentrar os raios nessas áreas. Matematicamente, se a PDF $p(\omega)$ for perfeitamente proporcional ao integrando $f(\omega)$, a variância cai para zero com uma única amostra (embora isso seja impossível na prática para cenas complexas).

\subsection{Múltipla Amostragem de Importância (MIS)}

Um problema clássico ocorre quando temos uma fonte de luz pequena e intensa junto a um material altamente especular. Amostrar apenas a luz gera ruído no material; amostrar apenas a BRDF gera ruído na luz. O \textbf{MIS}, introduzido por Eric Veach, combina os estimadores de ambas as estratégias usando pesos heurísticos baseados em suas respectivas PDFs:
\begin{equation}
    w_i(x) = \frac{p_i(x)^\beta}{\sum_j p_j(x)^\beta}
\end{equation}
Isso garante que o renderizador use a melhor estratégia para cada região da cena, resultando em uma convergência muito mais rápida.

\subsection{Roleta Russa e Caminhos Infinitos}

Caminhos de luz poderiam, em teoria, sofrer infinitos rebotes. Para evitar loops infinitos e desperdício de processamento em raios que já perderam quase toda a sua energia (baixo peso), usamos a \textbf{Roleta Russa}.
\begin{enumerate}
    \item Escolhemos uma probabilidade de terminação $q$ (baseada no albedo do material).
    \item Geramos um número aleatório $\xi \in [0, 1)$.
    \item Se $\xi < q$, o caminho é terminado.
    \item Se sobreviver, o peso do raio é escalado por $1/(1-q)$ para manter o estimador não-viesado.
\end{enumerate}

\section{Hardware Acceleration e o Futuro}

O Ray Tracing, historicamente restrito a supercomputadores e renderizadores offline para cinema, atingiu o mercado de consumo através de GPUs modernas com núcleos de aceleração dedicados (\textit{RT Cores}).

\subsection{Aceleração em Hardware}

Processadores gráficos contemporâneos implementam a travessia de BVH e o teste de interseção raio-triângulo diretamente em circuitos lógicos especializados. Isso permite que jogos executem milhões de raios por segundo, viabilizando iluminação global, reflexos dinâmicos e sombras fisicamente corretas em tempo real. A API \textbf{DirectX Raytracing (DXR)} e extensões do \textbf{Vulkan} padronizaram o acesso a esses recursos.

\subsection{Denoising e Inteligência Artificial}

Mesmo com hardware potente, gerar centenas de amostras por pixel em 16ms é impossível. A solução moderna é renderizar com pouquíssimas amostras (1 a 4 SPP) e utilizar filtros de \textbf{Denoising} (redução de ruído) espaciais e temporais. Técnicas baseadas em Deep Learning, como o NVIDIA DLSS, Intel XeSS e AMD FSR, utilizam vetores de movimento e dados históricos para reconstruir imagens limpas e de alta resolução a partir de entradas ruidosas e de baixa fidelidade.

\section*{Exercícios}

\begin{enumerate}
    \item Derive os coeficientes da equação quadrática para a interseção raio-esfera $\|\mathbf{O} + t\mathbf{D} - \mathbf{C}\|^2 = r^2$.
    \item No algoritmo de Möller-Trumbore, explique matematicamente o que ocorre quando o determinante é nulo. Qual o significado geométrico deste caso para a relação entre o raio e o plano do triângulo?
    \item Escreva o pseudocódigo para a travessia de uma árvore BVH usando uma pilha (\textit{stack}). Como você otimizaria a ordem de visita dos nós filhos para acelerar a busca?
    \item Demonstre matematicamente que a Roleta Russa mantém o estimador de Monte Carlo não-viesado (\textit{unbiased}). Por que o peso do raio deve ser dividido pela probabilidade de sobrevivência?
    \item Dada uma fonte de luz esférica de raio $R$ a uma distância $d$ do ponto $p$, como você realizaria \textit{importance sampling} diretamente sobre o ângulo sólido subtendido pela luz?
    \item Explique as diferenças físicas e as unidades de medida entre radiância, irradiância e fluxo radiante. Qual dessas grandezas é preservada ao longo de um raio no vácuo?
    \item Compare a complexidade computacional da rasterização vs. ray tracing em termos de número de primitivas ($N$) e número de pixels ($P$). Como as estruturas de aceleração mudam essa relação?
    \item O que é a aproximação de Schlick para o termo de Fresnel? Explique como ela varia com o ângulo de incidência e o índice de refração do material.
    \item Descreva o modelo de microfacetas GGX. Por que sua distribuição de normais é considerada superior ao modelo de Phong para materiais metálicos?
    \item Explique o conceito de Multiple Importance Sampling (MIS). Desenhe um cenário (configuração de luz e material) onde a amostragem da luz falha mas a da BRDF funciona, e vice-versa.
    \item Explique o fenômeno de \textit{total internal reflection} (Reflexão Total Interna) usando a Lei de Snell. Em que condições ele ocorre e como deve ser tratado no código de refração?
    \item O que é a SAH (\textit{Surface Area Heuristic})? Deduza a fórmula de custo e explique por que a área de superfície é um bom preditor da probabilidade de interseção de um raio aleatório.
\end{enumerate}




\begin{dica}
    No Path Tracing, o ruído diminui com a raiz quadrada do número de amostras ($1/\sqrt{N}$). Para reduzir o ruído visível pela metade, é necessário quadruplicar o número de amostras por pixel (SPP).
\end{dica}

\section*{Exercícios}

\begin{enumerate}
    \item Derive os coeficientes da equação quadrática para a interseção raio-esfera e resolva para $t$.
    \item No algoritmo de Möller-Trumbore, explique matematicamente o que ocorre quando o determinante é nulo. Qual o significado geométrico deste caso?
    \item Escreva o pseudocódigo para a travessia de uma árvore BVH usando uma pilha (\textit{stack}) de forma eficiente.
    \item Demonstre matematicamente que a Roleta Russa mantém o estimador de Monte Carlo não-viesado.
    \item Dada uma fonte de luz esférica, como você realizaria \textit{importance sampling} diretamente sobre a área da luz em vez de amostrar o hemisfério uniformemente?
    \item Explique as diferenças físicas entre radiância, irradiância e fluxo radiante.
    \item Compare a complexidade computacional da rasterização vs ray tracing conforme o número de triângulos e luzes na cena cresce.
    \item Como o fenômeno de \textit{total internal reflection} afeta o cálculo de refração em materiais vítreos? Use a Lei de Snell para justificar.
    \item O que é o modelo de microfacetas GGX e como ele difere do modelo de sombreamento de Phong?
    \item Explique o conceito de Multiple Importance Sampling (MIS) e dê um exemplo de cenário onde ele é essencial para reduzir a variância.
\end{enumerate}

